\clearpage\begingroup\SourceLeaf{16}\HistoricalHead{11}{R}{ÜBER DIE HANNOVERSCHEN NORMALPFUNDE.}
Ich hatte also die Wahl unter zwei Wegen, um das Hannoversche Normalpfund
darzustellen; entweder es so genau wie möglich dem preussischen Urpfunde
D an Masseninhalt gleich zu machen, oder so genau wie möglich
467711 Milliontheile der Grundeinheit A darzustellen.

Um zuvörderst den \textls[120]{ersten} Weg zu versuchen, hatte ich durch den Oberbergrath
\textsc{Schaffrinsky} in Berlin das Pfund E anfertigen lassen, welches vor
der Absendung von jenem und dem Professor \textsc{Encke} vermittelst wiederholter
Wägungen am 21. und 23. März 1836 mit dem preussischen Urpfunde D verglichen
[wurde]; das Resultat davon war, dass im Mittel E auf der Waage um
ein halbes Milligramm schwerer erschien als D. Obgleich die mir in extenso
mitgetheilten Protokolle auch den Barometer- und Thermometerstand während
der Wägungen enthielten, so fehlten doch nun noch die wesentlichern Elemente
behufs der Reduction auf den leeren Raum, nemlich die specifischen
Gewichte von D und E. Das letztere konnte ich nun freilich durch Abwägen
im Wasser selbst bestimmen, sobald mir die nöthigen Hülfsmittel zu Gebote
standen, und es ist diess auch später zu wiederholten mahlen geschehen.
Allein in Beziehung auf das preussische Urpfund D erfuhr ich nun erst mit
Bestimmtheit, dass dessen specifische Schwere gar nicht durch Versuche ausgemittelt,
also genau zu reden, dass sie \textls[120]{unbekannt} ist. Es folgt daraus, dass
unter diesen Umständen unmöglich ist, von D eine Copie zu machen, bei
welcher nicht eine Ungewissheit über die Gleichheit des wahren Inhalts zurückbleibt,
die bis auf 5 Milligramme gehen kann. Ich musste folglich entweder
auf eine grössere Schärfe Verzicht leisten, oder auf dem zweiten Wege zur
ersten Quelle zurückgehen.

Bei dem französischen Kilogramm A tritt nun zwar, wie schon oben bemerkt
ist, dieselbe ursprüngliche Unbestimmtheit des specifischen Gewichts ein: allein
diesem Uebelstande hat der Conferenzrath \textsc{Schumacher}, welcher sich mehr als
irgend ein anderer Gelehrte mit der schärfsten Vergleichung der verschiedenen
Normalgewichte beschäftigt hat, durch eine schätzbare Arbeit abgeholfen.
Sein eigenes Kilogramm B wurde in Paris, wohin zu dem Ende zweimahl
eine Reise unternommen ist, mit dem Originalkilogramm auf das sorgfältigste
verglichen, und zwar nicht bloss nach dem Gewichte auf der Waage, sondern
zugleich nach dem Raumvolumen\textsuperscript{*})\begingroup\renewcommand{\thefootnote}{}\footnotetext{\textsuperscript{*}) Die Anwendbarkeit dieses Verfahrens beruht auf dem Umstande, dass die Kilogramme A und B sehr nahe vollkommene Cylinder sind, und fällt ganz weg bei Gewichtsstücken, die wie das preussische Urpfund D keine so einfache Gestalt haben.}\endgroup\ vermittelst eines künstlichen besonders
\HistoricalFoot{}{2\textsuperscript{*}}
\clearpage\SourceLeaf{17}\HistoricalHead{12}{L}{PHYSIK, AMTLICHE BERICHTE.}
\noindent dazu verfertigten Apparats. Durch diese zum Theil in \textsc{Schumachers} Jahrbuch
für 1836 beschriebenen Arbeiten[\textsuperscript{*})]\begingroup\renewcommand{\thefootnote}{}\footnotetext{[\textsuperscript{*}) \textit{Vergleichung des Kilogramms von Platina, welches Etatsrath Schumacher aufbewahrt, mit dem gesetzlichen Kilogramm der Archive}, \textsc{Schumachers} Jahrbuch 1836, S. 237—250.]}\endgroup\ ist nun sowohl das bis dahin unbekannt
gewesene specifische Gewicht des französischen Normalkilogramms A, als das
Verhältniss von \textsc{Schumachers} Kilogramm B zu jenem auf das schärfste ausgemittelt.

Ich liess durch \textsc{Repsold} in Hamburg zwei Halbkilogramme C aus Messing
anfertigen, welche der Conferenzrath \textsc{Schumacher} die Gefälligkeit hatte, vor
der Absendung durch eine grosse Anzahl von Wägungen mit seinem Kilogramm
B zu vergleichen, und mir die Protokolle darüber in extenso mitzutheilen.

Es ist leicht zu erachten, dass die bei der Wahl des zweiten Weges zu
lösende Aufgabe, nemlich]

ein Gewichtsstück so zu reguliren, dass es nach seinem wahren Inhalte
\(\frac{467711}{1000000}\) von einem andern gegebenen werde, viel schwieriger und zeitraubender
ist, als die, wo es bloss darauf ankommt, ein Gewicht einem andern gleich
zu machen. Den grössten Zeitaufwand erfordern bei der ersten Aufgabe die
Vorbereitungen, insofern man erst sich selbst ein ganzes System von Hülfsgewichten
bilden und diese sämmtlich, vom grössten bis zum kleinsten herab,
mit einem viel grössern Grade von Schärfe bestimmen muss, als bei irgendwelchen
andern Gelegenheiten nöthig ist.

Meine auf jenen Zweck bezüglichen Arbeiten, die mich während des
Herbstes 1836 beschäftigt hatten, musste ich, ehe ich noch zu definitiven
Resultaten gelangt war, abbrechen, um zunächst meine ganze Zeit auf diejenigen
Arbeiten zu wenden, die sich auf die Darstellung der Hohlmaasse,
imgleichen auf die Darstellung der Medicinal- und Juweelengewichte bezogen.
Diese Umstände, nachher vielfache andere Arbeiten, endlich auch die Erwägung,
dass es angemessen sei, die Bildung der Normalgewichte bis zuletzt
ausstehen zu lassen, um dabei auch alle inzwischen eingesammelten Erfahrungen
benutzen zu können — war die Ursache, dass jene Arbeiten erst im Sommer
1839 wieder aufgenommen werden konnten. Es ergab sich dann aber bald
\clearpage\SourceLeaf{18}\HistoricalHead{13}{R}{ÜBER DIE HANNOVERSCHEN NORMALPFUNDE.}
\noindent das eben so unerwartete als unangenehme Resultat, dass die Halbkilogramme
C (obwohl stets unter Verschluss auf das sorgfältigste aufbewahrt und auch
im äussern Ansehen gar keine Veränderung zeigend) an ihrem Gewichtsinhalt
\textls[120]{Veränderungen} in den \(2\frac{1}{2}\) Jahren erlitten hatten, die zwar in jeder sonstigen
praktischen Beziehung als unerheblich gelten möchten, aber den hier
in Rede stehenden delicaten Zweck insofern vereitelten, als sie mich nöthigten,
die Arbeit vom Herbst 1836 ganz zu cassiren. Die Ursache jener Erscheinung
kann übrigens nur in dem Umstande gesucht werden, dass jene
Halbkilogramme C nicht vergoldet sind; eben deswegen aber konnte ich mich
nun auch auf die völlige Unveränderlichkeit des Systems von Hülfsgewichten,
welches ich 1836 gebildet hatte, nicht mehr verlassen, weil auch diese Hülfsgewichte
von Messing und unvergoldet waren, und so war ich gezwungen, mir
erst wieder ein neues Hülfssystem zu formiren.

Es bleibt mir nun noch übrig, die Art und Weise, wie ich durch die
vom August-November 1839 ausgeführten Arbeiten die Hannoverschen Normalpfunde
definitiv dargestellt habe, zu berichten. Das Kilogramm C hätte unter
den erwähnten Umständen, um als Zwischenglied zu dienen, erst von neuem
mit dem Kilogramm B verglichen werden müssen; anstatt dessen zog ich vor,
nunmehr das Pfund F als Zwischenglied zu gebrauchen. Für den Conferenzrath
\textsc{Schumacher} war die Arbeit ziemlich dieselbe, ob er C von neuem mit
B vergleichen, oder F in Bruchtheilen von C bestimmen sollte, da er im Besitz
eines von ihm selbst mit grösster Sorgfalt gebildeten Systems von Hülfsgewichten
ist, welches aus Platin verfertigt ist, also die möglich grösste Unveränderlichkeit
verbürgt. Ich selbst gewann aber dadurch den Vortheil, dass
ich das, wie eben erwähnt, neu gebildete System messingner Hülfsgewichte
nur behufs meiner Abwägungen im destillirten Wasser anzuwenden nöthig
hatte.

Zuerst wurde also vermittelst wiederholter Abwägungen in destillirtem
Wasser die specifische Schwere der verschiedenen hier zur Vergleichung kommenden
Gewichtsstücke bestimmt, und solche gefunden
\begin{center}
\begin{tabular}{@{}rlcl@{}}
für & E & .\quad.\quad. & 8,401109\\
    & F & .\quad.\quad. & 8,074584\\
    & G & .\quad.\quad. & 8,085009\\
    & H & .\quad.\quad. & 8,057088.
\end{tabular}
\end{center}

\clearpage\SourceLeaf{19}\HistoricalHead{14}{L}{PHYSIK, AMTLICHE BERICHTE.}
Hiernächst wurden diese vier Pfundstücke durch eine sehr grosse Anzahl
von Wägungen unter einander verglichen, und dann F nach Altona gesandt,
wo der Conferenzrath \textsc{Schumacher} dessen Gewichtsinhalt in Theilen des Kilogramms
B gleichfalls durch zahlreiche Wägungen auf das schärfste bestimmte.
Um nun aber alle Zweifel zu heben, ob nicht der Transport den Gewichtsinhalt
von F irgend wie verändern könne, wurde die Versendung mehreremale
wiederholt, und neue eben so zahlreiche Wägungen sowohl in Göttingen wie
in Altona angestellt; auf diese Weise hat das Pfund F die Reise zwischen
beiden Orten sechsmal gemacht und ist in drei verschiedenen Zeitperioden in
Göttingen mit E, G und H und eben so oft in Altona mit B verglichen.

Das Resultat aus 158 Wägungen in Göttingen war, dass (im leeren Raume)
\[
\left.\begin{array}{@{}lrl l@{}}
\text{E um}&0{,}025&\text{Milligramm}&\text{leichter}\\
\text{G um}&1{,}483&\text{»}&\text{leichter}\\
\text{H um}&2{,}989&\text{»}&\text{schwerer}
\end{array}\right\}\ \text{ist als F.}
\]
\noindent Dagegen war das Resultat der Wägungen in Altona, nach allen Reductionen,
dass das wahre Gewicht F 467706,310 Milliontheile des französischen Urkilogramms
A beträgt. Die wahren Gewichte von E, G, H sind folglich
\begin{center}
\begin{tabular}{@{}lrlr@{}}
E & 467706,285 & Milligramme & ,\\
G & 467704,827 & » & ,\\
H & 467709,299 & » & .
\end{tabular}
\end{center}
\noindent\hspace*{\parindent}Durch Hineinlegen sorgfältig abgeglichener Drahtstückchen, die respective
6,173 und 1,701 Milligramme betragen, sind endlich G und H auf den gesetzlichen
Inhalt von 467711 Milligrammen gebracht, so genau also normirt,
wie es heut zu Tage die Kunst vermag. Es wird natürlich vorausgesetzt, dass
künftig die Köpfe nicht wieder abgeschraubt werden. Um eine anschauliche
Vorstellung von der Kleinheit der letzten Correction zu geben, welche das
nach Hannover geschickte Pfund H erhalten hat, habe ich meinem Bericht
ein eben so grosses also 1,701 [Milligramm] schweres Drahtstückchen beigelegt

Es wird nicht ohne Interesse sein, hier noch anzuführen, was sich über
das Verhältniss der Hannoverschen beschriebenermaassen corrigirten Normalpfunde
zu dem preussischen Urpfunde D sagen lässt. Das Pfund E ist zwar,
nach seinem wahren Inhalt um 4,715 Milligramme leichter als G oder H
(nach der Correction durch die Drahtstücke): allein wegen seiner bedeutend
grössern specifischen oben angeführten Schwere verliert es beim Wägen in der
\clearpage\SourceLeaf{20}\HistoricalHead{15}{R}{ÜBER DIE HANNOVERSCHEN NORMALPFUNDE.}
\noindent Luft weniger als jene, und zwar im Durchschnitt aus allen vorgekommenen
Wägungen E um 4,027 Milligramme weniger als G und 4,397 Milligramme
weniger als H; oder mit andern Worten, bei einem mittleren Zustande der
Atmosphäre \textls[120]{erscheint} auf der Waage E um 0,688 Milligramm leichter als
G und um 0,318 Milligramm leichter als H. Da nun E um 0,5 Milligramm
schwerer auf der Waage erschien als D, so kann man, so weit man sich auf
die in Berlin gemachte Abwägung verlassen und absolute Unveränderlichkeit
von E seit fast 5 Jahren voraussetzen kann, schliessen, dass auf der Waage
G etwa um 1,2 und H etwa 0,8 Milligramm schwerer erscheinen würde,
als das Berliner Urpfund D. Also auch von dieser Seite ist die Uebereinstimmung
so gross, wie sie nur gewünscht oder erwartet werden dürfte.

\medskip
\noindent\hspace*{\parindent}Göttingen 27. Januar 1841.\hfill /: unterz. :/\quad C. F. \textsc{Gauss}.
\par\bigskip
\begin{flushright}\begin{minipage}{.66\textwidth}\small
\setlength{\parindent}{1.2em}
Die Uebereinstimmung dieser Abschrift mit dem Original wird hiermit bescheinigt.

\medskip Hannover, den 13. April 1841.
\begin{center}gez. \textsc{Bening}, Kanzleirath.\end{center}
\end{minipage}\end{flushright}
\bigskip
{\small
\begin{center}BEMERKUNG.\end{center}
\noindent\hspace*{\parindent}Zu den beiden vorstehenden Berichten vergleiche man den Briefwechsel zwischen \textsc{Gauss} und \textsc{Schumacher} vom 1. März 1836 bis zum 29. November 1839, siehe \textit{Briefwechsel zwischen Gauss und Schumacher} II, 1860, S. 439 bis III, 1861, S. 340, ferner \textsc{Schumachers} Abhandlung \textit{A Comparison of the late Imperial Standard Troy Pound weight with a Platina copy of the same, and with other standards of authority}, Philosophical Transactions of the Royal Soc. of London, 1836 II, S. 457—495.
\par\hfill\textsc{Schaefer}.
\par}
\begin{center}\rule{1.35in}{.35pt}\end{center}
\par\clearpage\endgroup
